\documentclass[11pt]{article}
\usepackage[margin=1in]{geometry}
\usepackage{amsmath,amssymb,amsthm,mathtools}
\usepackage[T1]{fontenc}
\usepackage{lmodern,microtype}
\usepackage[colorlinks=true,urlcolor=blue,linkcolor=blue]{hyperref}
\newtheorem{theorem}{Theorem}
\newtheorem{lemma}[theorem]{Lemma}
\theoremstyle{remark}\newtheorem{remark}[theorem]{Remark}
\newcommand{\ddc}{dd^c}
\newcommand{\avg}{\mathop{\mathrm{Avg}}}
\title{Uniform exact-adjoint endpoint bounds\\
on smooth complex-Hessian families}
\author{Research working paper}
\date{9 October 2026: independent verification pending}
\begin{document}
\maketitle
\begin{abstract}
On a fixed compact complex manifold, we give an explicit infinite
supremum-norm-bounded family of smooth logarithmic densities generated
by uniformly positive, bounded complex Hessians. The proposed theorem
bounds the determinant-weighted probabilities and determinant densities
of every exact maximizing current uniformly in $L^\infty$.
The proof controls the Gauduchon factor through its cofactor divergence:
the lower vector coefficient involves only the fixed background torsion.
We expand the required uniform Harnack argument. No derivative bound
on the varying Hessian or logarithmic density enters the final constants.
The family is nontrivial on every fixed manifold in the original
zero-lower-mass setting. This is an unverified author candidate for the
restricted estimate in original S17; unrestricted density families
and the original all-support problem P16 remain unresolved.
\end{abstract}

\section{The original setting and the explicit family}
Fix a connected compact complex manifold $X$ of dimension $n\geq3$,
a smooth positive Hermitian form $h$, and a smooth positive probability
volume $dV$. Throughout,
\[
 d^c=\frac{i}{4\pi}(\bar\partial-\partial),\qquad
 \ddc=\frac{i}{2\pi}\partial\bar\partial.
\]
Retain the original S17 assumption $v_-(h)=0$ for bounded genuine
$h$-psh potentials. Fix $0<\varepsilon_0<1$ and define
\[
 \mathcal V_{\varepsilon_0}=
 \{v\in C^\infty(X,\mathbb R):\ \sup_Xv=0,
             -\varepsilon_0h\leq\ddc v\leq\varepsilon_0h\}.
\]
For each $v$ put
\[
 \sigma_v=h+\ddc v,\quad J_v=\sigma_v^n/dV,\quad
 c_v=\int_XJ_v\,dV,\quad F_v=\log(J_v/c_v).
\]
These are actual smooth positive solution forms and smooth real
logarithmic densities, with $\int_Xe^{F_v}dV=1$.
Boundedness of the family $\{F_v\}$ uses the supremum norm on $C^0(X)$.
No derivative norm of $F_v$ or $\sigma_v$ is part of the family bounds.

For clarity, the only original maximizing-current background we retain is
as follows. Let $\mathcal G_h$ consist of positive $(n-1,n-1)$ currents
$T$ with $\ddc T=0$ and $\int_Xh\wedge T=1$. Original S17 defines
the determinant measure $\Gamma_F(T)$, an attained maximum
$M(F)=e^{-B_0(F)/n}$, and the weighted probability
$\nu_{T,F}=\Gamma_F(T)/M(F)$ at an exact maximizer. For a smooth $F$
the exact maximizing current is unique, smooth and strictly positive:
it is the normalized $(n-1)$st power of the Gauduchon metric conformal
to the smooth solution form. If $T=\gamma^{n-1}$ is that current, its
determinant density is
\[
 k_T=(\gamma^n/dV)^{(n-1)/n},\qquad
 \Gamma_F(T)=e^{F/n}k_T\,dV.
\]
These are the original problem's stated background, not conclusions
about arbitrary tests or about all supporting probabilities.

\begin{theorem}\label{thm:family}
There are constants $0<c<C<\infty$, depending only on
$X,h,dV,\varepsilon_0$, such that for every
$v\in\mathcal V_{\varepsilon_0}$ and every original exact maximizing
current $T$ for $F_v$, both densities obey
\[
 c\leq\frac{d\nu_{T,F_v}}{dV}\leq C,
 \qquad c\leq k_T\leq C.
\]
In particular, both have a uniform critical
$L^{n/(n-1)}(dV)$ bound. Moreover,
\[
 \|F_v-F_0\|_\infty
 \leq n\log\frac{1+\varepsilon_0}{1-\varepsilon_0},
\]
and the family contains infinitely many distinct normalized nonconstant
densities on the specified fixed zero-lower-mass manifold.
No new existence assertion for a zero-lower-mass manifold is made.
\end{theorem}

\section{Cofactors and the fixed torsion}
We first allow any smooth real $v$ with
$\delta h\leq\sigma=h+\ddc v\leq\Lambda h$.
Put $Q=\sigma^{n-1}$ and $J=\sigma^n/dV$.
Gauduchon's smooth conformal theorem
\cite[Th\'eor\`eme 2 and its full proof]{Gaud} gives a smooth
strictly positive function $\rho$ such that
\[
 \ddc(\rho Q)=0,\qquad\int_X\rho\,dV=1.
\]
The factor is unique with this normalization. The even first-Betti
condition in that article concerns its Theorem 1, and is not used.
Since $d\sigma=dh$, smooth differentiation gives
\begin{equation}\label{eq:dQ}
 dQ=(n-1)dh\wedge\sigma^{n-2}.
\end{equation}
Thus $Q$ and $dQ$ have bounded coefficients controlled only by
$\delta,\Lambda$ and fixed geometry. No derivative of the variable
complex Hessian occurs in this bound.

Define a real symmetric tensor $A$ and a vector $b$ intrinsically by
\begin{align*}
 \langle A\,d\phi,d\psi\rangle dV
       &=n\,d\psi\wedge d^c\phi\wedge Q,\\
 (b\cdot d\phi)dV&=-n\,d^c\phi\wedge dQ.
\end{align*}
Only the $(1,1)$ part of $d\psi\wedge d^c\phi$ contributes; it equals
\[
 \frac{i}{4\pi}
  (\partial\psi\wedge\bar\partial\phi
       +\partial\phi\wedge\bar\partial\psi).
\]
This proves symmetry and positivity of $A$.
It is uniformly elliptic and bounded in fixed real reference geometry.
Equation~\eqref{eq:dQ} bounds $b$ uniformly.
For $L_\sigma\phi=\operatorname{tr}_\sigma\ddc\phi$, put $P=JL_\sigma$.
Exact Stokes, with ordinary wedge powers, gives
\begin{align*}
 \int_X\psi P\phi\,dV
 &=-n\int_Xd\psi\wedge d^c\phi\wedge Q
        +n\int_X\psi d^c\phi\wedge dQ\\
 &=-\int_X\langle A\,d\phi,d\psi\rangle dV
        -\int_X\psi b\cdot d\phi\,dV.
\end{align*}
Indeed the last term of $d(\psi d^c\phi\wedge Q)$ is
$-\psi d^c\phi\wedge dQ$.
Consequently
\[
 P=\operatorname{div}_{dV}(A\,d)-b\cdot d,
 \qquad
 \operatorname{div}_{dV}(A\,d\rho+b\rho)=0.
\]
The latter follows from $\ddc(\rho Q)=0$. We need no sign or bound
on $\operatorname{div}b$, or on all component derivatives of $A$.
In a coordinate chart the fixed smooth volume density can be
absorbed into $A,b$, preserving uniformly controlled ellipticity
and the bound on the lower vector.

\section{A uniform Harnack argument}
\begin{lemma}\label{lem:Harnack}
Let $\rho$ be a smooth strictly positive solution on a fixed chart of
\[
 \operatorname{div}(A\nabla\rho+B\rho)=0,
 \qquad\lambda I\leq A\leq\lambda_1 I,\quad |B|\leq K.
\]
On a fixed-ratio inner cube,
$\sup\rho\leq C\inf\rho$, where $C$ depends only on
dimension, ellipticity, $K$ and the fixed chart size.
No derivative bound for $A,B$ enters this constant.
\end{lemma}
\begin{proof}
All powers below are admissible tests: the individual smooth
positive function has finite positive and negative moments on
each compact cutoff support. Their sizes are not assumed uniform.
For $p\ne0,1$ test by $\eta^2\rho^{p-1}$. Move the principal
energy to one side, using $|p-1|$ for either sign, and apply
Young's inequality to the principal cutoff, vector-gradient and
vector-cutoff terms. This gives
\begin{equation}\label{eq:power}
 \int\eta^2|\nabla\rho^{p/2}|^2
 \leq Cp^2(1+|p-1|^{-2})
          \int(|\nabla\eta|^2+K^2\eta^2)\rho^p.
\end{equation}
The coefficient bounds alone control $C$.
Write $m=2n$ and $\beta=m/(m-2)>1$.
Sobolev applied to $\eta\rho^{p/2}$ on nested cubes gives positive
power iteration for $p>0$, and the same iteration for $\rho^{-1}$
when $p<0$. Choose $p_0=\beta^{-k-1/2}$ with $k$ large.
Then $p_0\beta^j$ never equals one; the exponents near one have
a fixed dimensional gap. Geometric cutoff gaps and convergence
of $\sum_j\beta^{-j}(j+1)$ give
\begin{align}
 \sup_{Q/2}\rho&\leq C(\avg_Q\rho^{p_0})^{1/p_0},\label{eq:sup}\\
 \inf_{Q/2}\rho&\geq C^{-1}(\avg_Q\rho^{-p_0})^{-1/p_0}.
                                                        \label{eq:inf}
\end{align}
Polynomial growth of the constants in~\eqref{eq:power}, away
from $p=1$, is sufficient for these convergent products.

The test $\eta^2/\rho$ gives
\[
 \int\eta^2|\nabla\log\rho|^2
       \leq C\int(|\nabla\eta|^2+K^2\eta^2).
\]
For every subcube $S\subset Q$, use a cutoff on $2S$ inside a
larger fixed chart cube. Scaling and Poincare give
\[
 \avg_S|\log\rho-(\log\rho)_S|\leq K_0,
\]
with a uniform $K_0\geq1$.
Here and below $\avg$ denotes the Euclidean average.

We expand the exponential bridge. For any function $f$ with this
mean oscillation bound, fix $a>2$. Select maximal dyadic subcubes
$S$ of $Q$ where $\avg_S|f-f_Q|>aK_0$.
Their total measure is at most $|Q|/a$.
Their parents have average at most $aK_0$, hence
$|f_S-f_Q|\leq2^maK_0$.
Off their union differentiation gives $|f-f_Q|\leq aK_0$.
Repeat inside each selected cube with its own mean. Generation
$j$ has measure at most $a^{-j}|Q|$; off the next generation,
the accumulated increments and final oscillation are at most
$(j2^ma+a)K_0$. The intersection of all generations is null.
Integrating this exponential tail proves, for a dimensional $c>0$,
\[
 \avg_Q e^{c|f-f_Q|/K_0}\leq C_m.
\]
Apply it to $f=\log\rho$ and choose the earlier $k$ so that
$p_0<c/K_0$. It follows that
\[
 (\avg_Q\rho^{p_0})(\avg_Q\rho^{-p_0})\leq C.
\]
Combining with~\eqref{eq:sup}--\eqref{eq:inf} proves Harnack.
The selected $p_0$ and all constants depend only on the stated
uniform data, never on derivatives of the variable coefficients.
\end{proof}

A fixed finite chart cover and connectedness supply a finite
Harnack chain. Thus $\sup_X\rho\leq C_H\inf_X\rho$ with a
uniform $C_H$. Since $dV(X)=\int_X\rho\,dV=1$,
\begin{equation}\label{eq:rho}
 C_H^{-1}\leq\rho\leq C_H.
\end{equation}
This quantitative estimate supplements Gauduchon's qualitative
smooth conformal existence. Its constants depend only on fixed
geometry, $\delta,\Lambda$ and the fixed torsion $dh$.

\section{Exact-current and determinant normalization}
For a smooth $v$ define
\[
 M_0=\int_XJ\rho\,dV,\qquad
 T_* =\rho Q/M_0,\qquad r=J\rho/M_0.
\]
Smooth $\ddc$-Stokes with $\ddc(\rho Q)=0$ gives
\[
 \int_Xh\wedge\rho Q
  =\int_X\sigma\wedge\rho Q
  =\int_XJ\rho\,dV=M_0.
\]
Thus $T_*\in\mathcal G_h$. It is the $(n-1)$st power of the
smooth Gauduchon metric
$\gamma=(\rho/M_0)^{1/(n-1)}\sigma$ conformal to the actual
solution form. The original smooth exact-maximizer statement
identifies every exact maximizing current with $T_*$.
For $c_v=\int_XJdV$ and $F=\log(J/c_v)$ the actual solution has
$B_0(F)=\log c_v$, so the original scalar identity gives
$M(F)=c_v^{-1/n}$. Here $c_0$ is reserved for $v=0$.
The ordinary wedge normalization of its determinant yields
\[
 k_{T_*}=(\gamma^n/dV)^{(n-1)/n}
       =\frac{\rho}{M_0}J^{(n-1)/n}
       =J^{-1/n}r.
\]
Consequently
\[
 \Gamma_F(T_*)=c_v^{-1/n}r\,dV,
 \qquad \nu_{T_*,F}=r\,dV,\qquad\int_Xr\,dV=1.
\]
This calculation estimates the original exact determinant-weighted
probability, with its precise trace and wedge normalizations.
It asserts nothing about singular determinant measures of arbitrary
test currents, near-maximizers, or representation of all supports.

Let $J_h=h^n/dV$ and put
\[
 j_-=(1-\varepsilon_0)^n\min_XJ_h>0,\qquad
 j_+=(1+\varepsilon_0)^n\max_XJ_h.
\]
For our entire family, $j_-\leq J\leq j_+$ and
$j_-\leq M_0\leq j_+$ because $\int_X\rho dV=1$.
Equation~\eqref{eq:rho}, with
$\delta=1-\varepsilon_0$ and $\Lambda=1+\varepsilon_0$, gives
\[
 \frac{j_-}{C_Hj_+}\leq r\leq\frac{C_Hj_+}{j_-},
 \qquad
 j_+^{-1/n}\frac{j_-}{C_Hj_+}
       \leq k_{T_*}\leq
 j_-^{-1/n}\frac{C_Hj_+}{j_-}.
\]
These explicit bounds prove the density assertions of
Theorem~\ref{thm:family}. The critical norm is bounded by the
same upper constant since $dV$ is a probability volume.

\section{Boundedness and nontriviality at zero lower mass}
The inequalities
$(1-\varepsilon_0)^nJ_h\leq J_v\leq(1+\varepsilon_0)^nJ_h$
also hold after integration for $c_v/c_0$.
Subtracting the two logarithms gives
$\|F_v-F_0\|_\infty\leq
n\log((1+\varepsilon_0)/(1-\varepsilon_0))$.

Choose any smooth nonconstant real $\phi$ on the fixed $X$.
For all sufficiently small $t$, subtract the supremum of $t\phi$;
the resulting potential belongs to $\mathcal V_{\varepsilon_0}$.
The subtraction changes no Hessian or $F$.
Smooth differentiation at $t=0$ gives
\[
 \left.\frac{d}{dt}F_{t\phi}\right|_{t=0}
   =L_h\phi-
     \frac{\int_XJ_hL_h\phi\,dV}{\int_XJ_h\,dV}.
\]
This function is nonconstant. Otherwise $L_h\phi$ is constant;
evaluation at a maximum and a minimum makes that constant zero,
and the compact smooth strong maximum principle makes $\phi$
constant. The latter is the same smooth elliptic principle used
in Gauduchon's primary proof.
Choose two points where the displayed derivative differs. The
difference $F_{t\phi}(x)-F_{t\phi}(y)$ is strictly monotone on a
small parameter interval. It gives infinitely many distinct
normalized densities, with at most one parameter having zero
difference at these two points. Discarding that parameter yields
infinitely many nonconstant densities.
Distinct members cannot differ merely by a constant, because each
has $\int_Xe^FdV=1$.

The original metric $h$ and its hypothesis $v_-(h)=0$ have been
retained throughout. Restricting to this uniformly positive
smooth-Hessian family gives no positive lower mass for the full
bounded $h$-psh family. The construction uses the given fixed
zero-lower-mass manifold and does not require replacing it by a
K\"ahler manifold or constructing an additional such manifold.

\paragraph{Verification boundary.}
The original smooth maximizing-current and scalar-duality background
is used exactly as supplied by S17. Smooth Gauduchon conformal
existence and the classical smooth maximum principle are established
foundations. The cofactor identities, uniform weak estimates,
dyadic bridge, every-exact-current normalization and nontrivial
family calculations above are proposed proofs requiring independent
reconstruction and review. No pending sibling manuscript or prior
verdict is a correctness grant. No literature priority is asserted.
This source remains an unverified draft and is not published.

\begin{thebibliography}{9}
\bibitem{Gaud} P. Gauduchon,
\emph{Les m\'etriques standard d'une surface complexe compacte
\`a premier nombre de Betti pair}, Ast\'erisque 126 (1985), 129--135.
General Th\'eor\`eme 2, its positive smooth adjoint proof,
and the preceding smooth maximum principle.
\href{https://www.numdam.org/article/AST_1985__126__129_0.pdf}{Primary text}.
\end{thebibliography}
\end{document}
